\subsection{Neutrino Masses from Large Extra Dimensions --- arXiv:hep-ph/9811448}

\textbf{Authors:} Nima Arkani-Hamed, Savas Dimopoulos, Gia Dvali, John March-Russell

\paragraph{主な主張}
large extra dimension において右巻きニュートリノを bulk に置くことで体積希釈により自然に小さい Dirac 質量を生成でき、さらに遠方 brane での lepton number breaking を bulk messenger が伝えることで小さい Majorana 質量も生成できることを示した \cite{ArkaniHamed:1998LEDnu}。

\paragraph{新規性}
高い seesaw scale を導入せず、重力が弱く見えるのと同じ余剰次元の幾何学を用いてニュートリノ質量の小ささを説明する二つの高次元機構を具体化した点が新しい。

\paragraph{位置付け}
bulk right-handed neutrino とその KK tower を用いる large-extra-dimension neutrino phenomenology の原点の一つであり、後の active--sterile oscillation を通じた IceCube などの余剰次元探索へ直接つながる基礎論文である。
